\section{Cached option-score trajectories}

The broad panel uses cached layerwise scores for \ModelCount{} instruction-tuned open-weight checkpoints from the Qwen 2.5 Instruct, Gemma 2 IT, and Llama Instruct families \citep{qwen25report,gemma2report,llama3herd}. The benchmark slice contains \ArcPromptCount{} ARC-Challenge prompts \citep{clark2018arc} and \CsqaPromptCount{} CommonsenseQA prompts \citep{talmor2019commonsenseqa}. For each model, prompt, layer, and answer option, the cache stores the score assigned to the exact option text as the continuation after \texttt{Final answer:}. No auxiliary classifier is trained.

\begin{table*}[t]
\centering
\small

\begin{tabularx}{\linewidth}{@{}lX@{}}
\toprule
\textbf{Aspect} & \textbf{Release setting}\\
\midrule
Broad trajectory panel & Nine instruction-tuned open-weight checkpoints from Qwen 2.5 Instruct, Gemma 2 IT, and Llama Instruct on ARC-Challenge and CommonsenseQA. \\
Instrumented panel & Three 7--8B checkpoints on a 3,000-prompt MMLU cache with attention/MLP traces, span counterfactuals, and cached counterfactual accounting. \\
Primary readout & Exact-option continuation scores after the prompt suffix \texttt{Final answer:}. \\
Primary quantities & Leader, runner-up, entropy, soft margin, persistent contender depth $d_2$, persistent winner depth $d_1$, persistent 50\% final-gap depth $d_{50}$, late write, and trajectory regimes. \\
Main claim type & Cached readout-space dynamics plus mechanistic accounting of margin drift; no full circuit-discovery claim. \\
\bottomrule
\end{tabularx}

\caption{\textbf{Design summary.} The paper has a broad exact-option trajectory panel and a separate instrumented panel for mechanistic accounting.}
\label{tab:design}
\end{table*}

Let $w^\star$ be the final-layer winner, defined as the highest-scoring option at the final layer with deterministic tie handling. At layer $\ell$, let $\mathrm{rank}_\ell(w^\star)$ be the rank of this final winner among all answer options, where rank 1 is best. Let $g_\ell$ be the score gap between $w^\star$ and the strongest competing option at layer $\ell$, and let $g_L$ be the final-layer gap. When $g_L > 0$, define the normalized gap $m_\ell = g_\ell / g_L$.

We use three persistent depths:
\begin{align}
d_2 &= \min \{\ell : \mathrm{rank}_t(w^\star) \le 2 \text{ for all } t \ge \ell\}, \\
d_1 &= \min \{\ell : \mathrm{rank}_t(w^\star) = 1 \text{ for all } t \ge \ell\}, \\
d_{50} &= \min \{\ell : m_t \ge 0.50 \text{ for all } t \ge \ell\}.
\end{align}
Thus $d_2$ is persistent contender status, $d_1$ is persistent winner status, and $d_{50}$ is persistent decisive margin. The lags $\Delta_{21}=d_1-d_2$ and $\Delta_{250}=d_{50}-d_2$ measure how much depth separates candidate formation from winner isolation.

The mechanistic panel uses cached layerwise MMLU traces \citep{hendrycks2021mmlu} for three 7--8B instruction-tuned models. It contains margin drift, attention and MLP scalar contributions, prompt-span attention summaries, span-deletion effects, and cached counterfactual accounting. We use it only after the broad readout phenomenon is established.
